% =====================================================================
% Paper B — "Certified Explanation Reliability" — ICLR 2027 submission
% v1 DRAFT (extraction from the monograph per TWO_PAPER_SPLIT.md).
%
% TODO before submission (owner: Drake):
%   [ ] swap preamble for iclr2027_conference.sty when released (this
%       draft uses a compile-anywhere article preamble at ~ICLR density)
%   [ ] OSF DOI + CC-18 prospective results into Section 5 (placeholders
%       marked \osfpending)
%   [ ] anonymize artifact links (anonymous.4open.science) — ICLR is
%       double-blind; the Lean repo + monograph citations de-anonymize
%   [ ] co-author sign-off (Arnold, Rhoads) on the full claim set
% =====================================================================
\documentclass[10pt]{article}
\usepackage[margin=1in]{geometry}
\usepackage{amsmath,amssymb,amsthm}
\usepackage[numbers]{natbib}
\usepackage{booktabs}
\usepackage{microtype}
\usepackage{xcolor}
\usepackage{enumitem}
\usepackage[colorlinks=true,linkcolor=blue,citecolor=blue,urlcolor=blue]{hyperref}

% AUTO-GENERATED by paper/scripts/gen_claims.py — DO NOT EDIT.
% Generated 2026-07-31 from commits main:b514b8e dash:97a8451 ostrowski:080a55c
\newcommand{\ClaimMainFiles}{122}
\newcommand{\ClaimMainTheorems}{626}
\newcommand{\ClaimMainAxioms}{2}
\newcommand{\ClaimMainNativeDecide}{0}
\newcommand{\ClaimDashFiles}{59}
\newcommand{\ClaimDashTheorems}{368}
\newcommand{\ClaimDashAxioms}{2}
\newcommand{\ClaimDashNativeDecide}{0}
\newcommand{\ClaimOstrowskiFiles}{38}
\newcommand{\ClaimOstrowskiTheorems}{481}
\newcommand{\ClaimOstrowskiAxioms}{0}
\newcommand{\ClaimOstrowskiNativeDecide}{1}
\newcommand{\ClaimOverlapDashPct}{88}
\newcommand{\ClaimOverlapOstrowskiPct}{11}
\newcommand{\ClaimDistinctTheorems}{1090}
\newcommand{\ClaimTotalAxioms}{4}
\newcommand{\ClaimTotalNativeDecide}{1}


\newtheorem{theorem}{Theorem}
\newtheorem{proposition}{Proposition}
\newcommand{\snr}{\mathrm{SNR}}
\newcommand{\flip}{\mathrm{flip}}
\newcommand{\stab}{\textsc{stable}}
\newcommand{\marg}{\textsc{marginal}}
\newcommand{\unrel}{\textsc{unreliable}}
\newcommand{\osfpending}{\textcolor{red}{[OSF DOI + prospective results pending registration]}}

\title{Certified Explanation Reliability:\\ A Machine-Checked, Distribution-Free
Per-Query Certificate for Feature-Importance Claims}
\author{Anonymous authors\\ \small (double-blind submission)}
\date{}

\begin{document}
\maketitle

\begin{abstract}
Feature-importance explanations are unstable under model multiplicity: equally accurate
models routinely disagree about which feature matters more, so a published claim
``feature $j$ is more important than feature $k$'' may not survive retraining. Global
skepticism, however, is not actionable. We give a \emph{per-query} certificate that says
which individual ordering claims are reliable and which are not. For an ensemble of
equally performing models, the claim's signal-to-noise ratio ($\snr$) yields a
distribution-free guarantee: the fraction of models on which the claim flips is at most
$1/(1+\snr^2)$, by Cantelli's inequality applied to the empirical ensemble measure. The
guarantee, its \emph{tightness} (no certificate reading only the ensemble mean and
variance can promise less), its \emph{one-sidedness} (instability is provably not
certifiable from these statistics), and its \emph{transfer} to independent deployment
ensembles (a Hoeffding composition, plus a dependence-free variant with no independence
assumption at all) are machine-checked in Lean~4 --- to our knowledge the first empirical
ML instrument whose end-to-end guarantee chain is formally verified. Retrospectively, on
104 tabular datasets (10{,}724 ordering claims), every \stab{} verdict computed on one
ensemble stayed stable on an independent one (100\% transfer), the certificate's flip
ranking held out of sample (Spearman median $0.82$), its probability estimates were
calibrated (ECE $0.009$), and certified top-$k$ sets reproduced exactly in 100\% of cases
versus 27--30\% uncertified. A pre-registered prospective replication on a disjoint data
source (OpenML-CC18) is reported whatever its outcome \osfpending. Two case studies show
the stakes: on real credit data, 84\% of denied applicants' principal adverse-action
reasons flip across equally good models, while the certificate identifies the minority
of reasons that are dependable; on GPT-2, the certificate selects the ${\sim}10\%$ of
sparse-autoencoder features that reproduce across retraining ($0.81$ cosine
out-of-sample, versus $0.32$ uncertified). The certificate is ${\sim}50$ lines of code
around any ensemble.
\end{abstract}

\section{Introduction}
\label{sec:intro}

The same data and the same accuracy target admit many models --- the Rashomon effect
\citep{breiman2001statistical,fisher2019models,semenova2022existence,rudin2024amazing}
--- and those models disagree about their explanations. Post-hoc attribution methods
disagree with each other \citep{krishna2022disagreement}, are fragile to perturbation
\citep{ghorbani2019interpretation,slack2020fooling}, can fail basic sanity checks
\citep{adebayo2018sanity}, and, even for a fixed method, retraining an equally accurate
model can reverse the reported ranking \citep{caraker2026dash,caraker2026limits}. A
companion theory paper proves this is structural, not incidental: no explanation of an
underspecified system is simultaneously faithful, stable, and decisive
\citep{caraker2026limits}; related impossibility results constrain what attributions can
deliver in principle \citep{bilodeau2024impossibility,sundararajan2020many,kumar2020problems}.

An impossibility theorem, however, does not tell a practitioner what to do with the
explanation in front of them. Total distrust is as useless as total trust: some
explanation claims \emph{are} stable across every equally good model, and those are
exactly the ones safe to act on. The missing instrument is a \emph{per-query reliability
certificate}: for each individual claim, a guarantee --- not a heuristic score --- of how
likely that claim is to survive model multiplicity.

This paper supplies that instrument and validates it at scale. The certificate is
deliberately minimal. Retrain (or bootstrap) an ensemble of $M$ equally performing
models; for the ordering claim ``feature $j$ ranks above feature $k$,'' compute the
per-model importance differences $D_1,\dots,D_M$, their mean $\hat\mu$ and standard
deviation $\hat\sigma$, and the signal-to-noise ratio $\snr = |\hat\mu|/\hat\sigma$.
Then the fraction of models disagreeing with the claimed direction is at most
$1/(1+\snr^2)$ --- Cantelli's one-sided inequality applied to the empirical measure of
the ensemble, so the bound is \emph{distribution-free and assumption-free}: it holds for
every ensemble, exactly, with no appeal to Gaussianity, independence across features, or
asymptotics. An $\snr \ge 2$ (\stab{}) claim flips on at most $20\%$ of equally good
models, unconditionally.

\paragraph{Contributions.}
\begin{enumerate}[leftmargin=1.5em,itemsep=1pt,topsep=2pt]
\item \textbf{A machine-checked guarantee chain.} The finite-ensemble Cantelli bound
  (Theorem~\ref{thm:cantelli}), its tightness (Theorem~\ref{thm:tight}), its
  one-sidedness (Theorem~\ref{thm:onesided}), the family-wise simultaneous bound
  (Proposition~\ref{prop:simul}), and the out-of-sample transfer theorems
  (Theorems~\ref{thm:transfer-iid} and~\ref{thm:transfer-general}) are all formalized in
  Lean~4 \citep{demoura2021lean4} over Mathlib \citep{mathlib2020}, with no custom
  axioms. The paper's guarantee \emph{is} the verified statement; each theorem below
  names its Lean declaration.
\item \textbf{Transfer as a theorem, not an observation.} The bound is computed
  in-sample on one ensemble; deployment happens on another. We prove the guarantee
  transfers: with an explicit $e^{-2mt^2}$ rate for independent deployment ensembles,
  and --- removing the independence assumption entirely --- a population-Cantelli form
  valid for \emph{any} correlated or exchangeable ensemble. This closes the
  in-sample-to-out-of-sample gap that a validation study alone cannot.
\item \textbf{Validation at scale, in the correct epistemic mode.} We validate the
  \emph{bounds by coverage} (does the guarantee hold?), not by goodness-of-fit of any
  point-predictive law. On 104 PMLB datasets and 10{,}724 ordering claims: 100\% \stab{}
  transfer to independent ensembles, out-of-sample flip ranking at Spearman median
  $0.82$, calibration ECE $0.009$, and certified top-$k$ reproduction 100\% versus
  27--30\% uncertified (Section~\ref{sec:validation}).
\item \textbf{A pre-registered prospective round.} Thresholds for five confirmatory
  predictions on a disjoint benchmark suite (OpenML-CC18) were frozen and registered
  before any CC-18 data was accessed; the analysis pipeline is frozen at a public commit
  hash. Results are reported whatever the outcome (Section~\ref{sec:prospective}).
\item \textbf{Two high-stakes case studies.} Legally mandated adverse-action reasons on
  real credit data (84\% flip across the Rashomon set; the certificate finds the
  dependable minority), and sparse-autoencoder features on GPT-2 (the certificate's
  \stab{} calls transfer at $0.81$ cosine out-of-sample versus $0.32$ uncertified)
  (Section~\ref{sec:cases}).
\end{enumerate}

Equally important is what we do \emph{not} claim; Section~\ref{sec:scope} states the
scope honestly, including the certificate's provable inability to certify
\emph{instability}, its per-model-class nature, and the falsified point-predictive
extensions of the underlying theory that this paper deliberately does not rest on.

\section{The certificate and its verified guarantees}
\label{sec:certificate}

\paragraph{Setting.} Fix a dataset and a model class, and let $\mathcal{E} = \{m_1,
\dots, m_M\}$ be an ensemble of models with (near-)equal predictive performance ---
in our experiments, bootstrap-retrained random forests; in general, any Rashomon-set
sample \citep{semenova2022existence}. An importance functional $I$ (impurity, SHAP
\citep{lundberg2017unified}, permutation, \dots) assigns each model a vector
$I(m)\in\mathbb{R}^p$. For an ordered pair $(j,k)$, the \emph{claim} is
$\mathrm{sign}(\mathbb{E}[D])$ where $D_i = I(m_i)_j - I(m_i)_k$; the \emph{flip
fraction} is the share of models whose $D_i$ disagrees with the claimed direction. The
certificate computes
\[
\hat\mu = \tfrac1M\textstyle\sum_i D_i,\qquad
\hat\sigma^2 = \tfrac1M\textstyle\sum_i (D_i-\hat\mu)^2,\qquad
\snr = |\hat\mu|/\hat\sigma,
\]
and reports the band \stab{} ($\snr\ge 2$), \marg{} ($\snr \ge 0.5$), or \unrel{},
together with the guarantee of Theorem~\ref{thm:cantelli} and the calibrated estimate
$\Phi(-\snr)$.

\begin{theorem}[Distribution-free flip bound; Lean:
\texttt{cantelli\_lower\_tail}, \texttt{flip\_le\_one\_div\_one\_add\_snrSq} in
\texttt{CertificateGuarantee.lean}]
\label{thm:cantelli}
For any ensemble with $\hat\mu \neq 0$,
\[
\flip(\mathcal{E})\;\le\;\frac{1}{1+\snr^2}.
\]
In particular a \stab{} claim ($\snr\ge2$) flips on at most $20\%$ of the ensemble
(Lean: \texttt{stable\_flip\_le\_one\_fifth}).
\end{theorem}

\emph{Proof idea.} A flip is the one-sided event $D - \hat\mu \le -|\hat\mu|$ under the
empirical measure of the ensemble; Cantelli's inequality bounds it by
$\hat\sigma^2/(\hat\sigma^2 + \hat\mu^2)$. Because the inequality is applied to the
empirical measure --- the finite list of retrained models the certificate actually sees
--- there are no distributional hypotheses to discharge: the bound holds for every
observed ensemble, always. It is strictly sharper than two-sided Chebyshev
($1/\snr^2$) and never vacuous. \hfill$\square$

Two further verified facts pin down exactly what such a certificate can and cannot
promise.

\begin{theorem}[Tightness; Lean: \texttt{flip\_bound\_tight} in
\texttt{CertificateTight.lean}]
\label{thm:tight}
For every rational tail fraction there is an explicit two-point ensemble attaining
$\flip = \hat\sigma^2/(\hat\sigma^2+\hat\mu^2)$ with equality. Consequently no
certificate that reads only the ensemble mean and variance can promise a smaller flip
rate: the $20\%$ \stab{} guarantee is exactly met, not conservative-by-design.
\end{theorem}

\begin{theorem}[One-sidedness; Lean: \texttt{flip\_zero\_at\_low\_snr}]
\label{thm:onesided}
There is an explicit ensemble with $\snr^2 < 1$ (the \unrel{} band) whose flip fraction
is exactly $0$. Hence low $\snr$ does not --- and provably cannot --- certify that a
claim is unstable: from $(\hat\mu,\hat\sigma)$ alone, stability is certifiable and
instability is not.
\end{theorem}

Theorem~\ref{thm:onesided} began as an attempt to prove the opposite (a
Paley--Zygmund-style lower band that would ``certify arbitrariness''); the counterexample
is the honest outcome, and it dictates the instrument's interface: the certificate
reports \stab{} or \emph{not-certified}, never ``certified unstable.'' This matters for
the regulatory use case in Section~\ref{sec:cases}.

\begin{proposition}[Simultaneous claims and certified top-$k$; Lean:
\texttt{simultaneous\_cantelli}]
\label{prop:simul}
For a family of claims, the fraction of the ensemble on which \emph{any} claim flips is
at most the sum of the per-claim Cantelli bounds. In particular, call a top-$k$ feature
set \emph{certified} when the boundary comparison ($k$-th versus $(k{+}1)$-th by mean
importance) is \stab{}; the union bound then controls the whole set.
\end{proposition}

\section{Transfer: the in-sample bound is an out-of-sample guarantee}
\label{sec:transfer}

Theorem~\ref{thm:cantelli} is a statement about the ensemble on which it is computed.
Deployment poses a different question: a practitioner certifies a claim on ensemble $A$
today and retrains ensemble $B$ tomorrow. The natural referee objection --- \emph{your
bound is in-sample; your evidence is out-of-sample; where is the link?} --- is answered
by two theorems, not by the experiments.

\begin{theorem}[Transfer, independent ensembles; Lean:
\texttt{transfer\_flipFreq\_bound}, \texttt{transfer\_flip\_bound} in
\texttt{TransferTheorem.lean}]
\label{thm:transfer-iid}
Model the deployment ensemble as $m$ i.i.d.\ draws with population flip probability $p
\le \beta$, where $\beta$ is the population Cantelli bound $1/(1+\snr_{\mathrm{pop}}^2)$.
Then for every $t>0$,
\[
\Pr\!\left[\widehat{\flip}_B \ge \beta + t\right]\;\le\;e^{-2mt^2}.
\]
\end{theorem}

The proof composes Mathlib's Hoeffding inequality with the population Cantelli bound ---
itself formalized as \texttt{cantelli\_upper} in \texttt{PopulationCantelli.lean}, filling
a gap in formalized probability libraries. In practice, however, retrained ensembles are
bootstrap-\emph{correlated}, not independent. Rather than assume the correlation away, we
remove the independence hypothesis entirely:

\begin{theorem}[Transfer, dependence-free; Lean: \texttt{transfer\_flip\_general} in
\texttt{TransferGeneral.lean}]
\label{thm:transfer-general}
For \emph{any} deployment ensemble --- arbitrarily correlated or exchangeable ---
population Cantelli applied to the flip frequency itself gives
\[
\Pr\!\left[\flip_B \ge \beta + t\right]\;\le\;
\frac{\operatorname{Var}(\flip_B)}{\operatorname{Var}(\flip_B) + t^2}.
\]
\end{theorem}

The bound degrades gracefully as correlation inflates the flip-frequency variance, with
Theorem~\ref{thm:transfer-iid} as the sharper independent-case idealization. Under
either, the 100\% observed transfer of Section~\ref{sec:validation} is a theorem's
predicted consequence rather than a fortunate regularity. To our knowledge this makes
the certificate the first empirical-ML instrument with an end-to-end formally verified
chain from guarantee statement to deployment behavior; formal verification in ML has
otherwise targeted models and robustness properties rather than empirical instruments
\citep{decker2024provably,noguer2025mathematical}.

\section{Retrospective validation at scale}
\label{sec:validation}

\paragraph{Protocol.} 104 PMLB classification and regression datasets ($200\le n\le
5000$, $4\le p\le 30$ after filtering), $M=60$ bootstrap-retrained random forests per
dataset with impurity importances, split 30/30 into ensemble $A$ (computes the
certificate) and an independent ensemble $B$ (measures the flip); all
$\binom{p}{2}$-style ordering claims, 10{,}724 in total. The validation mode is
\emph{coverage of bounds}: we ask whether guarantees hold, not whether a fitted curve
fits.

\paragraph{Results.}
\begin{itemize}[leftmargin=1.5em,itemsep=1pt,topsep=2pt]
\item \textbf{The guarantee holds.} The uncapped Cantelli bound holds for all 10{,}724
  claims in-sample, as the theorem requires. The composite out-of-sample bound
  (population Cantelli $+$ Hoeffding slack, Theorem~\ref{thm:transfer-iid}) covers
  $99.4\%$ of claims against its $95\%$ design level, and the dependence-free composite
  covers $100\%$. The bound is honestly conservative: mean slack ${\approx}0.50$,
  expected from Cantelli looseness plus the finite-$m$ Hoeffding term.
\item \textbf{\stab{} transfers perfectly.} Every \stab{} verdict computed on $A$ stayed
  stable (flip $\le 20\%$) on the independent $B$ --- $100\%$ of 4{,}521 \stab{} claims
  across all 100 datasets that produced \stab{} calls. \unrel{} claims flipped at median
  $0.42$, near coin-flip: the bands separate.
\item \textbf{The ranking is predictive out of sample.} Per-dataset Spearman between
  predicted flip $\Phi(-\snr_A)$ and observed flip on $B$: median $0.82$ (IQR
  $[0.75,0.88]$), beating the group-free $|\mathrm{corr}(X_j,X_k)|$ baseline on $98\%$
  of datasets and a single-model importance-gap baseline on $87\%$. The single-model gap
  is itself a competitive \emph{ranking} heuristic --- but returns no calibrated
  probability and no guarantee.
\item \textbf{Calibrated.} Pooled reliability diagram of $\Phi(-\snr_A)$ against
  observed flips: ECE $0.009$.
\item \textbf{Certified top-$k$ reproduces.} Across the 104 datasets, certified top-5
  sets (Proposition~\ref{prop:simul}) reproduced \emph{exactly} on the independent
  ensemble in $100\%$ of cases, versus $27$--$30\%$ for uncertified sets. Certification
  fires rarely --- few datasets have a stable top-$k$ boundary --- which is itself the
  point: most ``top-feature'' claims are not stable, and the certificate says precisely
  which are.
\item \textbf{Failures, named.} Two datasets (balance\_scale, xd6) show weak
  out-of-sample ranking; we report them rather than average them away.
\end{itemize}

\section{Pre-registered prospective validation (OpenML-CC18)}
\label{sec:prospective}

Retrospective validation, however large, cannot rule out pipeline-level overfitting to
the benchmark pool. We therefore froze the exact pipeline of
Section~\ref{sec:validation} (only the data loader, an NA-handling rule, ordinal
encoding, and a mechanical PMLB-deduplication step differ, each pre-specified) at a
public commit hash, and registered five confirmatory predictions on OSF \emph{before any
CC-18 data was downloaded}, with a mandatory waiting period between registration and
first download. The predictions (P1--P5, thresholds frozen in the registration; full
table in Appendix~\ref{app:prereg}) cover \stab{} transfer ($\ge 95\%$ per dataset in
$\ge 90\%$ of eligible datasets), pooled \stab{} flip magnitude, ranking value over
baselines, calibration, and certified top-$k$ reproduction, each with an explicit
disconfirmation criterion. \textbf{The paper reports the outcome whatever it is}: a
failed prediction narrows the claim's scope in the final version rather than
disappearing. \osfpending

\section{Case studies: what certification changes in practice}
\label{sec:cases}

\subsection{Adverse-action reasons on real credit data}
\label{sec:credit}

U.S.\ lenders must state the principal reasons for adverse credit decisions (ECOA /
Regulation~B; model-risk guidance \citep{occ2011sr117}), and analogous explanation
rights exist in the EU \citep{euaiact2024}; in practice reasons are derived from feature
attributions on a single model \citep{selbst2018intuitive,ustun2019actionable}. On
German credit data (OpenML-31), we trained 40-model Rashomon ensembles under a strict
$2\%$ equal-accuracy window (5 seeds) and computed each denied applicant's principal
(top-SHAP) reason per model. \textbf{On average $84\%$ of denied applicants (range
$76$--$93\%$ across seeds) have a principal reason that flips across equally good
models}: the legally operative explanation is, for most applicants, an artifact of which
equally accurate model was deployed. The certificate identifies the dependable
minority: it certifies a reason for ${\sim}19\%$ of denied applicants (range
$4$--$29\%$), and certified reasons are more concentrated across the ensemble (modal
reason share $0.85$ versus $0.71$ uncertified). By Theorem~\ref{thm:onesided} we state
the negative half carefully: the $84\%$ is an \emph{observed} instability across the
constructed Rashomon set, not a certified impossibility --- the certificate can prove
the minority stable but provably cannot prove the majority arbitrary. The compliance
upshot is constructive: report certified reasons where they exist; where none exists,
the honest notice is that no single principal reason is stable across equally accurate
models.

\subsection{Certified sparse-autoencoder features on GPT-2}
\label{sec:sae}

Sparse autoencoders are a leading tool for extracting interpretable features from
language models, but SAE features are themselves subject to multiplicity: retrained SAEs
learn overlapping-but-different dictionaries \citep{paulo2025sae,meloux2025everything}.
We trained an ensemble of 6 SAEs (2048 latents, $R^2\approx0.54$, $L_0\approx41$) on
GPT-2-small layer-6 residual activations and applied the certificate to per-feature
stability, evaluating against a held-out SAE ensemble. Features only partially reproduce
overall (mean best-match cosine $0.36$; random-direction null $0.12$), reproducing the
SAE-Rashomon finding at real-LLM scale. The certificate separates, out of sample: the
$9.5\%$ of features it certifies transfer to the held-out ensemble at mean cosine
$0.81$ ($88\%$ above $0.7$), while uncertified features sit near the ensemble average
($0.32$). The certificate thus turns ``SAE features don't fully escape multiplicity''
into a usable filter for which mechanistic-interpretability claims are worth trusting.
This study is exploratory --- one model, one layer, one SAE architecture --- and is
scoped as such.

\section{Honest scope}
\label{sec:scope}

\begin{itemize}[leftmargin=1.5em,itemsep=1pt,topsep=2pt]
\item \textbf{Per-model-class, not model-invariant.} A claim stable for one model class
  need not be stable for another: \stab{}-call agreement across gradient-boosting,
  random-forest, and linear ensembles is ${\approx}0.65$, and between-class instability
  exceeds within-class (median spread $0.058$ vs $0.032$). A certificate that pools three
  tree-ensemble classes does \emph{not} transfer better to a held-out linear model
  ($0.60$ vs $0.63$): cross-architecture reliability requires architecturally diverse
  agreement sets. All guarantees here are per ensemble-generating process.
\item \textbf{One-sided by theorem.} The certificate can certify stability, never
  instability (Theorem~\ref{thm:onesided}). ``Not certified'' means \emph{not proven
  stable}, nothing stronger.
\item \textbf{Conservative by design.} Cantelli tightness (Theorem~\ref{thm:tight}) is
  worst-case; on benign data the bound's mean out-of-sample slack is ${\approx}0.5$, and
  certification fires on a minority of claims (e.g.\ ${\sim}19\%$ of credit applicants,
  $9.5\%$ of SAE features). The instrument prefers silence to false confidence.
\item \textbf{Validation envelope.} The headline study is random forests with impurity
  importances on $n\gg p$ tabular data ($200\le n\le5000$, $4\le p\le30$); the SAE study
  is exploratory. No claim is made outside this envelope in this paper.
\item \textbf{What the underlying theory does not support.} Point-predictive extensions
  of the companion theory --- a quantitative law predicting instability magnitude from
  symmetry structure, and recovery of the exact stable-dimension from trained ensembles
  --- were tested and \emph{falsified} in prospective rounds reported in
  \citep{caraker2026limits}. This paper's claims are bounds validated by coverage, and
  none rests on the falsified arms.
\item \textbf{Known weak spots.} Two of 104 datasets show weak out-of-sample ranking
  (Section~\ref{sec:validation}); the credit certificate's strict-window separation is
  directional, not dispositive, at the achieved sample sizes
  (Section~\ref{sec:credit}).
\end{itemize}

\section{Related work}
\label{sec:related}

\textbf{Rashomon and predictive multiplicity.} The model-multiplicity literature
quantifies disagreement among equally good models
\citep{breiman2001statistical,fisher2019models,semenova2022existence,marx2024uncertainty,rudin2024amazing};
explanation disagreement specifically is documented by
\citet{krishna2022disagreement}. We add the per-claim guarantee layer: not \emph{whether}
models disagree, but \emph{which individual claims} survive disagreement, with a proof.
\textbf{Statistical inference for explanations.} Confidence intervals and tests for
SHAP-style attributions
\citep{watson2023testing,zhang2026statistical,jin2025probabilistic,hwang2025shap} and
stability selection \citep{meinshausen2010stability} provide per-quantity uncertainty;
our certificate differs in being distribution-free at the ensemble level, tight and
one-sided by theorem, and machine-checked end to end, including its out-of-sample
transfer. \textbf{Impossibility and critique.} Attribution impossibility and fragility
results \citep{bilodeau2024impossibility,sundararajan2020many,kumar2020problems,%
ghorbani2019interpretation,adebayo2018sanity,slack2020fooling} establish the ceiling;
the companion theory \citep{caraker2026limits,caraker2026dash} derives it from model
multiplicity and identifies the stable (symmetry-invariant) subspace. The certificate is
the constructive complement: it identifies, per query, the claims already inside the
stable region. \textbf{Verified ML.} Formal verification in ML has focused on model
properties (robustness, fairness) \citep{decker2024provably,noguer2025mathematical};
verifying the \emph{guarantee chain of an empirical instrument} appears new.

\section{The practitioner interface}
\label{sec:practice}

The certificate is deliberately small --- no new estimator, no tuning, ${\sim}50$ lines
around any ensemble:

\begin{quote}\small\ttfamily
fit M equally-performing models (bootstrap or re-seed);\\
D[i] = importance(model i, feature j) - importance(model i, feature k);\\
snr = |mean(D)| / std(D);\\
band = STABLE if snr >= 2 else MARGINAL if snr >= 0.5 else UNRELIABLE;\\
guarantee: flip <= 1/(1+snr\^{}2);\quad estimate: Phi(-snr);\\
top-k certified iff boundary pair (k, k+1) is STABLE.
\end{quote}

Reference implementation, the 104-dataset study, the frozen prospective runner, and the
Lean development are in the artifact (Appendix~\ref{app:artifact}).

\section{Conclusion}

Model multiplicity makes most feature-importance claims less solid than they look, and a
theory-level impossibility says this cannot be engineered away. The practical response is
not despair but triage: a per-query certificate with a distribution-free, tight,
provably one-sided, provably transferring guarantee --- every link machine-checked ---
that turns ``explanations are unreliable'' into ``\emph{these} claims are reliable,
\emph{those} are not certified.'' At scale, its promises held every time they were made.
The invitation to the community is adversarial: the pre-registered round is designed so
that failure would be visible, and the certificate's own theory says exactly what would
falsify it.

\bibliographystyle{plainnat}
\bibliography{references}

\appendix

\section{Pre-registered predictions (P1--P5)}
\label{app:prereg}

Thresholds frozen at registration; ``eligible'' means $\ge5$ \stab{} claims on ensemble
$A$. \osfpending

\begin{center}\small
\begin{tabular}{@{}lll@{}}
\toprule
 & Prediction & Threshold \\
\midrule
P1 & \stab{} transfer (flip $\le0.20$ on $B$) per dataset & $\ge95\%$, in $\ge90\%$ of eligible datasets \\
P2 & pooled mean \stab{} flip on $B$ & $\le0.05$; no dataset $>0.20$ \\
P3 & ranking: median Spearman; beats $|\mathrm{corr}|$ baseline & $\ge0.60$; on $\ge85\%$ of datasets \\
P4 & pooled calibration ECE & $\le0.05$ \\
P5 & certified top-5 exact reproduction (if $\ge10$ instances) & $\ge90\%$ certified vs $<60\%$ uncertified \\
\bottomrule
\end{tabular}
\end{center}

Disconfirmation criteria are stated verbatim in the registration: P1 failing falsifies
the deployability of the \stab{} band under data-source shift; P3/P4 failing falsifies
the quantitative (Gaussian-calibrated) arm independently of the band-level guarantee.
Both would be reported as such.

\section{Artifact and verification statement}
\label{app:artifact}

The Lean~4 development ({\ClaimMainFiles} files, {\ClaimMainTheorems} theorems across
the program's main repository) builds with zero \texttt{sorry}; the certificate-chain
theorems cited in this paper (Theorems~\ref{thm:cantelli}--\ref{thm:transfer-general},
Proposition~\ref{prop:simul}) depend only on Lean's core axioms
(\texttt{propext}, \texttt{Classical.choice}, \texttt{Quot.sound}), verified by
\texttt{\#print axioms} under continuous integration; the repository-wide count is
{\ClaimTotalAxioms} domain axioms, none on the results used here. The experiment scripts
reproduce Section~\ref{sec:validation} from public PMLB data with fixed seeds; the
prospective runner is frozen at the commit hash cited in the OSF registration.
% TODO(anonymization): replace repo pointers with anonymous.4open.science mirror for review.

\end{document}
